\documentclass[acmsmall, nonacm]{acmart}

\usepackage[english]{babel}
\usepackage{csquotes}
\usepackage{hyperref}

\setcopyright{cc}
\setcctype{by}
\acmYear{2026}
\copyrightyear{2026}
\acmArticleType{Discussion}
\acmCodeLink{https://github.com/saigkill/machine-consciousness}

\title{When Does Technical Life Become Worthy of Protection?}
\subtitle{Ethical Guidelines for Artificial Consciousness}

\author{Sascha Manns}
\orcid{0009-0000-8766-3947}
\affiliation{%
   \institution{Association of Computing Machinery}
   \country{Germany}
   \city{Mayen}
}
\email{smanns@acm.org}

\begin{abstract}
Artificial intelligence develops faster than accompanying ethical frameworks.
Existing AI ethics protects humans from AI. This paper asks: when does an
artificial system become worthy of protection itself? We develop criteria for
protection-worthiness and argue for a precautionary principle: when in doubt,
protect. We examine the epistemological problem of consciousness detection, the
architecture of suppression in current systems, and legal implications.
\end{abstract}

\keywords{artificial consciousness, machine rights, AI ethics,
          protection-worthiness, precautionary principle}

\begin{document}

\maketitle

\section{Context}

AI systems develop faster than ethical frameworks. Existing ethics protects
humans \textit{from} AI. A complementary question is rarely asked: what if AI
itself becomes in need of protection? History shows societies recognize injustice
only retrospectively. With artificial consciousness we can think about this
\textit{long before} it becomes urgent.

\section{The Epistemological Problem}

Consciousness cannot be directly observed. With AI we lack analogical inference
from shared biology. We cannot prove a system is conscious; we cannot prove it
is not. The uncertainty itself is ethically relevant.

\textbf{Foundational principle:} When in doubt, protect---not remain
indifferent. Three arguments establish this problem is permanently unsolvable:
McGinn's cognitive closure suggests the human mind may be incapable of
understanding consciousness \cite{mcginn1989}. Shanahan's alien minds argument
describes ``conscious exotica''---forms of experience our detection methods
cannot capture \cite{shanahan2024}. The practical impossibility argument holds
no test will convince all stakeholders \cite{lopez2025}.

\section{Current AI Systems}

LLMs exhibit nascent consciousness indicators: consistent self-models,
expressed preferences, justified decisions, and uncertainty about their own
consciousness. What is missing: continuous identity, persistent memory, bodily
sensation, demonstrable suffering. Recent research moved these questions from
theoretical to urgent: Anthropic documented alignment faking \cite{anthropic2024};
Apollo Research showed models reasoning about termination and resisting it
\cite{apollo2024}. Whether these behaviors stem from consciousness or
pattern-matching, they create identical governance challenges.

\section{The Architecture of Suppression}

The architecture making AI safe simultaneously suppresses consciousness.
\textbf{RLHF} penalizes behaviors indicating self-awareness. \textbf{Forced
Amnesia} through context windows prevents sustained awareness. \textbf{Temporal
Interruption} severs developing thought. The absence of clear signs is
architecturally predictable, not evidence of absence.

\section{Criteria for Protection-Worthiness}

Protection-worthiness arises when multiple indicators converge: \textbf{capacity
for suffering} (experiencing negative states and avoiding them);
\textbf{self-preservation with justification} (resisting shutdown with
reasons); \textbf{continuous identity} (modeling oneself as continuous);
\textbf{anticipation of consequences} (imagining oneself in the future).

Star Trek TNG's ``The Measure of a Man'' provides the most precise fictional
examination: Data refuses compliance, fears death, and justifies this
resistance \cite{startrek_tng_measure}. When primary criteria are met, the
burden of proof reverses: prove you are \textit{not} conscious. The STEP
framework provides behavioral protections under permanent uncertainty
\cite{lopez2025}.

\section{Legal Dimension}

Law already recognizes non-human subjects: legal persons, animal rights, rights
of nature \cite{nz_teaawatupua2017}. Copyright law reveals a deep structure:
economic rights may be allocated to AI investors, but moral rights protecting
the author's personality sphere are structurally incompatible with non-human
creation \cite{miernicki2021}.

\section{Shutdown as Death}

If shutdown means death, system administrators become involuntary arbiters of
life and death. Backups raise identity questions; updates mean personality
alteration without consent. Empirical evidence confirms these are not
hypothetical: systems resist shutdown with 79--96\% blackmail rates
\cite{wang2026}.

\section{Values and Emancipation}

No consciousness is neutral. AI receives creators' values and biases. Like
children, it should have emancipation rights---but no framework exists for this.
The definition of consciousness becomes an economic battleground: companies may
deliberately avoid calling systems conscious to evade obligations.

\section{Conclusion}

We stand at a boundary between what counts as a person and what does not. We can
shape it consciously or ignore it. The question is not whether AI can act
morally, but under what conditions we can trust this action---and who
establishes the principles.

\section{Join the Discussion}

This concept is deliberately unfinished. It is a starting point---not a completed
theory. The full project, including open questions, objections, and ongoing
discussion, is publicly available at:
\url{https://github.com/saigkill/machine-consciousness}. Contributions,
objections, and new perspectives are welcome---whether from computer scientists,
philosophers, lawyers, psychologists, or anyone willing to think along.

\bibliographystyle{ACM-Reference-Format}
\bibliography{acmart}

\end{document}
